% \subsection{Power Sets}
\begin{definition}
    [Power Set] If A is a set, the power set of A is another set, denoted as $p(A)$ and defined to be the set of all subsets of A. 
    In symbols, $p(A) = \left\{X  \mid  X \subseteq A\right\}$
\end{definition}
\begin{remark}
    If A is a finite set, then $\left|p(A)\right| = 2^{\left|A\right|}$
\end{remark}
% \subsection{Intersection, union, set difference}
\begin{definition}
    [Intersection, union, and set difference]
    Given two sets A and B, we define the following:
    \begin{itemize}
        \item Intersection $A \cap B = \left\{x \mid x \in A \text{ and } x \in B\right\}$
        \item Union $A \cup B = \left\{x \mid x \in A \text{ or } x \in B\right\}$
        \item Difference $A - B = \left\{x \mid x \in A \text{ and } x \notin B\right\}$
    \end{itemize}
\end{definition}

Observe that for any sets $X$ and $Y$ it is always true that $X \cup Y = Y \cup X$ and $X \cap Y = Y \cap X$, but in general $X - Y \neq Y - X$.

\begin{definition}
    [Complement]
    Let A be a set with Universal set U. The complement of A, denoted $\bar{A}$, is the set $\bar{A} = U - A$.
\end{definition}

When a mathematical problem involves lots of sets, it is often convenient to keep track of them by using subscripts.
Thus instead of denoting three sets $A$, B, and C, we might instead write them as $A_1$, $A_2$, and $A_3$. These are called indexed sets. 
\begin{definition}
    Supposed $A_1, A_2, \cdots, A_n$ are sets. Then 
    \begin{equation*}
        A_1 \cup A_2 \cup A_3 \cup \cdots \cup A_n
        =
        \left\{
        x \mid x \in A_i \text{ for at least one set } A_i,\ 1 \leq i \leq n
        \right\}
    \end{equation*}  
    \begin{equation*}
        A_1 \cap A_2 \cap A_3 \cap \cdots \cup A_n
        =
        \left\{
            x \mid x \in A_i \text{ for every set } A_i, \text{ for } 1 \leq i \leq n
        \right\}
    \end{equation*}
\end{definition}

Here are the notations for indexed sets:
Given sets $A_1, A_2, A_3, \cdots, A_n$, we define 
\begin{equation*}
    \bigcup_{i = 1}^{n} A_i = A_1 \cup A_2 \cup A_3 \cup \cdots \cup A_n
\end{equation*}
and 
\begin{equation*}
    \bigcap_{i = 1}^{n} A_i = A_1 \cap A_2 \cap A_3 \cap \cdots \cap A_n
\end{equation*}

The notation is also used when the list of sets $A_1, A_2, A_3, \cdots$ is infinite. 
\begin{equation*}
    \bigcup_{i = 1}^{\infty} A_i = A_1 \cup A_2 \cup A_3 \cup \cdots \cup A_n
    =
    \left\{
        x \mid x \in A_i \text{ for at least one set of } A_i, \text{ with } 1 \leq i
    \right\}
\end{equation*}
and 
\begin{equation*}
    \bigcap_{i = 1}^{\infty} A_i = A_1 \cap A_2 \cap A_3 \cap \cdots \cap A_n
     =
    \left\{
        x \mid x \in A_i \text{ for every set } A_i, \text{ with } 1 \leq i
    \right\}
\end{equation*}

We can also have another set of notation for indexed sets. 
\begin{definition}
    \[
        \bigcup_{a \in I} A_a = \left\{
            x \mid x \in A_a \text{ for at least one set } A_a \text{ with } a \in I
        \right\}
    \]
    \[
        \bigcap_{a \in I} A_a = \left\{
            x \mid x \in A_a \text{ for every set } A_a \text{ with } a \in I
        \right\}
    \]
\end{definition}